\section{The Key Arrives as a String}
\label{sec:ch06-the-key-arrives-as-a-string}

\index{reference resolver!what it is handed|(}%
\code{\_entities} is the other new field, and behind it is the method
section~\ref{sec:ch06-three-edits} named and did not write. Read how little of
the work it is before writing it. The specification gives the whole resolver as
three steps: make an array, walk the representations, return the array. The walk
has three substeps of its own, and exactly one of
them is yours~\autocite{apollosubgraphspec}:

\begin{quote}
Pass the full representation object to whatever mechanism the library provides
the subgraph developer for fetching entities of the corresponding
\code{\_\_typename}.
\end{quote}

\index{reference resolver!is one step and nothing else}%
The substeps on either side of it read the \code{\_\_typename} off the
representation and put what came back into the array, and the second of those
carries the rule chapter~\ref{ch:federation-model} already quoted: the objects
must be listed in the order their representations were. So the loop, the
ordering, the array and the dispatch on type name all belong to the package.
Apollo states the division rather than leaving it to be inferred: the subgraph
library is responsible \enquote{for providing the mechanism that developers use
to specify this logic, and for automatically hooking into this
mechanism}~\autocite{apollosubgraphspec}. The method you write answers one
representation and knows nothing about the others.

\subsection{Not the Key You Expect}

\index{reference resolver!receives the undecoded node id}%
Here is the thing that catches everyone who has read
chapter~\ref{ch:schema-design} first. The key field is \code{id}, and since
chapter~\ref{ch:schema-design} every \code{id} in this schema is a global node
id: the type name, a colon and the primary key, base64 encoded. The router sends
that string, and the reference resolver is handed the string.

\begin{minted}{json}
{"__typename": "Speaker", "id": "U3BlYWtlcjox"}
\end{minted}

\index{node resolver!is handed the decoded key}%
The node resolver next door is handed the integer~1. That decoding happens in
front of it, which is why chapter~\ref{ch:schema-design} could say the method
never sees an id it has to parse. On the entity route nothing decodes anything,
and the parameter that arrives is \code{U3BlYWtlcjox}.

\index{ID@\code{[ID<T>]}!does not decode here}%
The attribute that decodes an id everywhere else in this book does not decode
this one. Write the parameter as \code{[ID<Speaker>] int id} and the file
compiles, the schema is unchanged, and the resolver runs with \code{id} set
to~0. Nothing throws. The entity comes back null, which reads exactly like a key
that matched no row, and the only sign that anything is wrong is that it
happens for every key. That is a bad failure to debug and a worse one to deploy,
so take the string.

\index{INodeIdSerializer@\code{INodeIdSerializer}!does the decoding}%
Decoding it is one call on \code{INodeIdSerializer}, which you ask for by naming
it in the parameter list like any other service. Its \code{Parse} method takes
the string and the type of the key inside it, and the second argument is the one
to get right: it is the type of the key, not the type of the entity.
\code{typeof(int)}, because the primary key is an \code{int}. Pass
\code{typeof(Session)} instead and the call throws \code{The node ID string has
an invalid format.}, which is the message chapter~\ref{ch:schema-design} met
from a malformed id and which here means something else entirely.

\index{NodeId@\code{NodeId}!carries the type name and the key}%
What comes back has two parts. \code{nodeId.InternalId} is the integer, which is
what the resolver came for. \code{nodeId.TypeName} is the string
\code{Speaker}, read out of the same base64. The resolvers below spend the first
and ignore the second, which is the ordinary thing to do and which
section~\ref{sec:ch06-a-key-you-did-not-produce} is about.

\subsection{The Projection Does Not Come With You}

\index{QueryContext@\code{QueryContext<T>}!cannot be injected here}%
One more parameter does not work, and it is the one every other resolver in this
book takes. A \code{QueryContext} cannot be injected into a reference resolver.
Adding it compiles, the schema builds, the service starts, and the first
\code{\_entities} request comes back as an unexpected execution error with no
exception in any log to explain it. I am describing that response rather than
printing it, because a service in that state is not a state this book ships and
nothing here can reproduce it for you.

\index{projection!is absent from the entity route}%
So the entity route runs without the projection
chapter~\ref{ch:data-without-n-plus-1} put in, and
section~\ref{sec:ch06-what-the-entity-route-costs} measures the difference. The
resolver below does not have one, and the comment in it says why rather than
leaving a later reader to wonder whether it was forgotten.

\index{key@\code{[Key]}!on the record instead is silent}%
One more placement is worth knowing before you copy the attributes onto a type
of your own. \code{[Key("id")]} on the \code{Speaker} record rather than on
\code{SpeakerType} also compiles, and produces a schema with no key in it at
all. The attribute is read off the class the source generator is building the
type from, and a positional record is not that class. The union in
section~\ref{sec:ch06-three-edits} is where you would notice.

\subsection{Both Resolvers}

\index{SpeakerType@\code{SpeakerType}!gains a reference resolver}%
\code{SpeakerType.cs} in full, with everything chapter~\ref{ch:schema-design}
did not have marked. The reference resolver borrows the same DataLoader the node
resolver above it uses, which is a one-line decision with a measurable
consequence in the next section:

\begin{minted}[highlightlines={1-3,8-10,23-41}]{csharp}
using HotChocolate.ApolloFederation.Resolvers;
using HotChocolate.ApolloFederation.Types;
using HotChocolate.Types.Relay;

namespace Sessions;

[ObjectType<Speaker>]
[Key("id")]
[ReferenceResolver(
    EntityResolver = nameof(ResolveSpeakerReferenceAsync))]
public static partial class SpeakerType
{
    /// <summary>
    /// Loads one speaker by the identifier inside a global node id.
    /// </summary>
    [NodeResolver]
    public static async Task<Speaker?> GetSpeakerAsync(
        int id,
        ISpeakerByIdDataLoader speakerById,
        CancellationToken ct)
        => await speakerById.LoadAsync(id, ct);

    /// <summary>
    /// Answers one representation from the entity route. The key
    /// arrives as the undecoded node id string, where the node
    /// resolver above is handed the integer already. Going through the
    /// same DataLoader is what makes a batch of representations cost
    /// one statement rather than one each.
    /// </summary>
    [GraphQLIgnore]
    public static async Task<Speaker?> ResolveSpeakerReferenceAsync(
        string id,
        INodeIdSerializer serializer,
        ISpeakerByIdDataLoader speakerById,
        CancellationToken ct)
    {
        var nodeId = serializer.Parse(id, typeof(int));

        return await speakerById.LoadAsync(
            (int)nodeId.InternalId!, ct);
    }
}
\end{minted}

\index{GraphQLIgnore@\code{[GraphQLIgnore]}!keeps the resolver out of the schema}%
\code{[GraphQLIgnore]} is load-bearing. Without it the source generator sees a
public method on a type class and does what it does with public methods on type
classes: publishes it as a field. \code{Speaker} would gain a field taking an id
and returning a different \code{Speaker}, which is not a thing anyone meant to
put in an API.

\index{SessionType@\code{SessionType}!gains a reference resolver}%
\code{SessionType.cs} takes the same shape and differs in the two ways this
section has already named. It goes to the table directly rather than through a
loader, because that is what its node resolver does, and it has no
\code{QueryContext} because it cannot have one:

\begin{minted}[highlightlines={2-4,10-12,38-59}]{csharp}
using GreenDonut.Data;
using HotChocolate.ApolloFederation.Resolvers;
using HotChocolate.ApolloFederation.Types;
using HotChocolate.Types.Relay;
using Microsoft.EntityFrameworkCore;

namespace Sessions;

[ObjectType<Session>]
[Key("id")]
[ReferenceResolver(
    EntityResolver = nameof(ResolveSessionReferenceAsync))]
public static partial class SessionType
{
    /// <summary>
    /// The person giving this session.
    /// </summary>
    public static async Task<Speaker?> GetSpeakerAsync(
        [Parent(nameof(Session.SpeakerId))] Session session,
        ISpeakerByIdDataLoader speakerById,
        CancellationToken ct)
        => await speakerById.LoadAsync(session.SpeakerId, ct);

    /// <summary>
    /// Loads one session by the identifier inside a global node id.
    /// </summary>
    [NodeResolver]
    public static async Task<Session?> GetSessionAsync(
        int id,
        ConferenceContext db,
        QueryContext<Session> query,
        CancellationToken ct)
        => await db.Sessions
            .Where(s => s.Id == id)
            .With(query)
            .FirstOrDefaultAsync(ct);

    /// <summary>
    /// Answers one representation from the entity route. Two things
    /// differ from the node resolver above, and neither is a choice.
    /// The key arrives as the undecoded node id string, so this method
    /// decodes it. And a QueryContext cannot be injected here, so there
    /// is no projection and every column is selected. One statement per
    /// representation is what chapter 10 measures and then fixes.
    /// </summary>
    [GraphQLIgnore]
    public static async Task<Session?> ResolveSessionReferenceAsync(
        string id,
        INodeIdSerializer serializer,
        ConferenceContext db,
        CancellationToken ct)
    {
        var nodeId = serializer.Parse(id, typeof(int));
        var key = (int)nodeId.InternalId!;

        return await db.Sessions
            .Where(s => s.Id == key)
            .FirstOrDefaultAsync(ct);
    }
}
\end{minted}

\index{entities@\code{\_entities}!answering}%
That compiles, and the entity route answers. The operation is one line and the
chapter uses it several times, so it stands on its own here:

\begin{minted}{graphql}
query E($r: [_Any!]!) {
  _entities(representations: $r) {
    __typename
    ... on Session { title }
  }
}
\end{minted}

\index{entities@\code{\_entities}!on the wire}%
On the wire, with the operation text standing in for the block above and the
body wrapped between tokens to fit the page:

\begin{minted}{text}
POST /graphql HTTP/1.1
Host: localhost:5001
Content-Type: application/json

{"query":"<the operation above>","variables":
{"r":[{"__typename":"Session","id":"U2Vzc2lvbjox"}]}}
\end{minted}

\begin{minted}[fontsize=\footnotesize]{json}
{"data":{"_entities":[{"__typename":"Session","title":"Schemas That Outlive Their Authors"}]}}
\end{minted}

\index{entities@\code{\_entities}!order is the only correspondence}%
Send three representations and three objects come back in that order. Mix the
types and they still come back in that order, which is what the union is for and
what the substep after yours requires. Send a key that matches no row and that
position is \code{null} while its neighbors are unaffected, which is the
specification permitting exactly one thing to go wrong per representation.
Chapter~\ref{ch:representation-order} is about what a subgraph that gets the
ordering wrong looks like from outside, and the short version is that it looks
fine.

\begin{onhc14}
The package exists on Hot Chocolate~14 and restores without complaint, and
\code{AddApolloFederation()} is there and does something. \code{\_service},
\code{\_Service}, \code{\_Any}, \code{FieldSet}, the \code{@link} line and
\code{@shareable} on the paging type all arrive exactly as they do on~16.

What does not arrive is the entity half. \code{[Key]} and
\code{[ReferenceResolver]} on the type class compile without a warning and
reach the schema not at all: no \code{@key} on any type, no \code{\_entities}
on \code{Query}, no \code{\_Entity} union. The 14 source generator has no
mechanism for applying an arbitrary descriptor attribute declared on a type
class, so it never reads them. The failure is silent in the worst place: a
composer asking \code{\_service} gets a subgraph it cannot route into, rather
than an error.

It is also not inert. The method \code{[ReferenceResolver]} names is published
as a field, \code{[GraphQLIgnore]} and all, so \code{Speaker} gains
\code{resolveSpeakerReference(id: String!): Speaker} with the doc comment as
its description.

A federated entity on 14 has to leave the generator. Declare the type as a
code-first \code{ObjectType<Speaker>} and both directives apply, because the
descriptor calls are made rather than discovered. This is the whole of its
\code{Configure}:

\begin{minted}[fontsize=\footnotesize]{csharp}
protected override void Configure(
    IObjectTypeDescriptor<Speaker> descriptor)
{
    descriptor
        .ImplementsNode()
        .ResolveNodeWith(
            typeof(SpeakerType).GetMethod(
                nameof(ResolveByKeyAsync))!);

    descriptor
        .Key("id")
        .ResolveReferenceWith(
            _ => ResolveReferenceAsync(default!, default!));
}
\end{minted}

The two methods it points at decode the key exactly as the ones above do: the
same undecoded string arrives and \code{Parse} is the same call. Only the
attachment differs. The cost is the class itself, because
\code{[ObjectType<Speaker>]} and \code{ObjectType<Speaker>} are two
declarations of one type and collide, so a federated entity on 14 cannot be
written in the idiom the rest of this book uses. One further trap sits in the
listing: \code{ResolveNodeWith} takes a \code{MethodInfo} where
\code{ResolveReferenceWith} two lines below it takes an expression.
Appendix~\ref{app:hc14} carries the class in full.
\end{onhc14}
\index{reference resolver!what it is handed|)}%
